\begin{abstract}
Automatic prompt and LLM-pipeline optimizers (DSPy-style compilers, best-arm-identification and minibatch searches) are routinely run on a benchmark's own questions, and the tuned system's score on those same questions is reported. We study what that number means. We prove that the optimism of the post-search score is at least the search regret and at most $\sqrt{\ln K/2n}$, and equals $\sigma_n\sqrt{1-\rho}\,m_K$ for equally good configurations with outcome correlation $\rho$. We then measure it on real LLM outcomes (Qwen2.5-0.5B/1.5B-Instruct and SmolLM2-1.7B-Instruct on MMLU-Pro; five search spaces of 40--80 prompt/format/demonstration configurations; five optimizer types; 2000 episodes per setting), where true accuracy is known by construction. With a 100-question dev set the reported accuracy exceeds the true accuracy of the selected configuration by 3.9--5.3 points (13--33\% of the accuracy) and exceeds even the best configuration's true accuracy in 74--90\% of episodes; the gap shrinks only like $n^{-0.6}$, grows with search effort much faster than the selected configuration improves, and is predicted within 10\% by a Gaussian plug-in. It distorts conclusions: a paired test declares a gain over an already-strong baseline in 14--19\% of episodes (nominal 5\%), reported lifts of 4.3--7.0 points over a mid-ranked baseline correspond to true lifts of 0.3--2.0 points, and a system searched over 32 configurations looks 3--10$\times$ better against one searched over 4 than it is. In a harness that tests eight reporting protocols against known truth, sample splitting is exactly unbiased but noisy and costs deployed accuracy; a search-replay bootstrap removes 70--80\% of the optimism for any optimizer; empirical-Bayes shrinkage is most accurate when configurations are exchangeable but fails with clustered quality levels; max-$t$ intervals are always valid but conservative. We give recommendations and discuss limitations (cached discrete search spaces, small models).
\end{abstract}
